% The decode path, the generated source, and where a real encoder would plug in.
% Columns are on a fixed pitch of 2.2 cm so that a three-line block cannot touch
% its neighbour.
\begin{tikzpicture}[font=\sffamily\small, x=1cm, y=1cm]

% ---------------- the generated source
\node[modelled, text width=32mm] (gen)  at (-5.9,3.0)
  {generator timer\\{\scriptsize two compare outputs, a quarter period apart}};
\node[modelled, text width=32mm] (genz) at (-5.9,0.8)
  {index output\\{\scriptsize one pulse every N counts}};
\begin{scope}[on background layer]
\node[layer, fit=(gen)(genz), inner sep=6pt,
      label={[layerlabel, anchor=south west]north west:the stand-in source}] {};
\end{scope}

% ---------------- the decode path
\node[hw, text width=34mm] (iface) at (0.4,3.0)
  {encoder interface\\{\scriptsize both edges of both inputs: four counts per cycle}};
\node[hw, text width=34mm] (cap)   at (0.4,0.8)
  {capture channel\\{\scriptsize the count at the index edge, in hardware}};
\node[fw, text width=34mm] (ext)   at (0.4,-1.4)
  {extend to 64 bit, signed\\{\scriptsize the direction bit decides which way a wrap went}};
\begin{scope}[on background layer]
\node[layer, fit=(iface)(cap)(ext), inner sep=6pt,
      label={[layerlabel, anchor=south west]north west:real, at full rate}] {};
\end{scope}

% ---------------- what comes out
\node[user, text width=30mm] (pos)   at (6.2,3.0) {position, in counts\\{\scriptsize absolute after the first index}};
\node[user, text width=30mm] (vel)   at (6.2,0.8) {velocity\\{\scriptsize two methods, and the node says which}};
\node[user, text width=30mm] (state) at (6.2,-1.4) {node state\\{\scriptsize chapter 1's one structure}};

% ---------------- what is absent, and the harness
\node[absentblk, text width=32mm] (enc) at (-5.9,-2.6)
  {a rotary encoder\\{\scriptsize not on this bench. It would drive the same
   three pins and nothing above would change}};
\node[ext, text width=32mm] (harn) at (0.4,-4.2)
  {test harness\\{\scriptsize knows both counts, asserts they are equal}};

% ---------------- edges
\draw[flow] (gen)  -- node[lbl, above]{A, B} (iface);
\draw[flow] (genz) -- node[lbl, above]{Z} (cap);
\draw[flow] (iface) -- (ext);
\draw[flow] (cap)   -- (ext);
\draw[flow] (ext)   -- (state);
\draw[flow] (iface) -- (pos);
\draw[flow] (pos)   -- (vel);
\draw[flow] (vel)   -- (state);
\draw[hiz, ->] (enc.east) -- ++(1.1,0) |- node[lbl, above, pos=0.3]{the day one is bought} (cap.west);
\draw[flow, draw=linecol!60] (harn.west) -- ++(-1.9,0) |- node[lbl, above, pos=0.25]{drives} (genz.south);
\draw[flow, draw=linecol!60] (ext) -- node[lbl, right]{reads} (harn);

\node[umlnote, text width=52mm, anchor=north west] at (-9.6,-4.4)
  {The decoder has no reference to the generator, no shared state with it and no
   way to ask it anything. That is what lets a real encoder replace the
   generator later without the decoder changing at all, and it is what makes the
   equality test meaningful rather than circular: the only thing the two modules
   share is a wire.};
\end{tikzpicture}
